\documentclass[11pt]{article}

\usepackage[a4paper,margin=1in]{geometry}
\usepackage{booktabs}
\usepackage{xurl}
\usepackage{array}
\usepackage{microtype}
\usepackage{fancyhdr}
\usepackage[hidelinks]{hyperref}

\hypersetup{
  pdftitle={Model Intervention Records: Addendum on Weight Patches and Record Qualifiers},
  pdfauthor={MIR-1 contributors},
  pdfsubject={MIR-1: additions after version 1.0},
  pdfkeywords={weight-space programming, model intervention records, weight
    patches, conformance testing, reproducibility, evaluation validity}
}

\setlength{\parskip}{0.35em}
\setlength{\parindent}{0pt}

\pagestyle{fancy}
\fancyhf{}
\fancyhead[L]{\small\textsc{mir-1 addendum}}
\fancyhead[R]{\small Weight Patches and Record Qualifiers}
\fancyfoot[C]{\small\thepage}
\renewcommand{\headrulewidth}{0.4pt}
\renewcommand{\footrulewidth}{0pt}

\begin{document}

\thispagestyle{plain}
\begin{center}
\vspace*{0.4em}
{\LARGE\bfseries Model Intervention Records}\\[0.35em]
{\LARGE\bfseries Addendum: Weight Patches and Record Qualifiers}\\[0.8em]
{\large Addendum to version 1.0}\\[0.25em]
{\large 18 November 2026}\\[0.8em]
{\small\textbf{Keywords:} weight-space programming; weight patches; model
intervention records; conformance testing; reproducibility}
\end{center}
\vspace{0.2em}
\hrule
\vspace{0.6em}

\begin{abstract}
Version 1.0 of this work specified a record format and eleven checks. Two
documents were subsequently published beside it: one defines a composable
payload for the \texttt{weights} intervention type and five checks reported in a
separate block; the other names four measurements the evidence already carries.
This addendum records both in one citable place and states exactly what did and
did not change. No requirement was withdrawn, no member was added to a record,
no canonical byte changed, and no identifier moved. The eleven checks of version
1.0 are the eleven checks of suite 1.7.0. What was added is a wider payload and
a finer vocabulary for evidence, not a new test.
\end{abstract}

\section{Scope, and what ``additive'' means here}

The paper specifies form, not desirability: a record states an executable
falsifier, an invariant set with at least one metric outside the claim's own
scoring function, and content-addressed lineage to a base checkpoint, and a
mechanical test reports eleven verdicts. Two additions followed publication.
Neither is a revision, and nothing in the paper is withdrawn.

``Additive'' is used here as a measurement and not as a claim of intent. The
test is: a record that carries none of the additions must produce the report it
produced before. That was checked against the validator as it stood before the
additions (suite 1.5.0) on all six records in the release. Every check line ---
identifier, verdict, and detail string --- is identical, and every record
identifier is unchanged. What differs is the last line: a record that carries a
patch reports the patch block beside the eleven, and a record that earns a
qualifier gains one line naming it. The suite number moved; nothing else did.

\begin{table}[h]
\centering
\begin{tabular}{@{}ll@{}}
\toprule
Property & State after the additions \\
\midrule
Record schema & byte-identical (\texttt{schema.json}, 6{,}369 bytes) \\
Worked example & \texttt{MIR-R1}, identifier unchanged \\
Requirements & eleven, unchanged; no new required field \\
Suite version & 1.7.0 \\
\bottomrule
\end{tabular}
\end{table}

\section{Weight Patches}

A \texttt{weights} payload may now name a patch instead of a full checkpoint. A
patch states the base checkpoint by its hash, an ordered list of operations, any
patch it depends on, and an inverse. The operation vocabulary is closed:
\texttt{set}, \texttt{add}, \texttt{scale}, and \texttt{add\_vector}. Indices are
ascending and unique within an operation, so that two serializers cannot spell
one patch two ways, and the patch identifier is the hash of the canonical bytes
with the identifier field removed.

Composition is executable rather than described. A patch whose base is a
\texttt{wpp-array-0.1} container can be applied to the container, and a stack of
patches is applied in dependency order, which is a depth-first post-order of the
lineage graph. This is what lets the test do something other than read the
patch: it applies it, applies the inverse, and compares bytes.

Reversibility is reported in two forms because one of them is not always
available. Restoring a value by adding the negative of a delta is not
bit-identical in binary32: $(x+d)-d$ can differ in the last place for some $x$
and $d$. An \emph{exact} inverse therefore stores the prior values and restores
them by assignment, and only an exact inverse is executed. An
\emph{algebraic} inverse --- negation, or a reciprocal for a scale --- is
recorded for a patch that has no captured prior, and is reported as not run. A
test that executed it would fail additive patches for a reason that is
arithmetic rather than a defect in the patch.

Five checks, P-01 to P-05, report on a patch: the base resolves and matches the
record, the patch conforms and its identifier matches its bytes, dependencies
resolve in order with no overlapping writes, the patch restates the record's
invariant set and falsifier, and the round trip is exact. They are reported in a
separate block and are never merged into the eleven. That separation is
deliberate: the eleven are the format's requirements, the five are a property of
one payload, and a record without a patch must produce the report it produced
before. Three qualifiers --- \texttt{patch-executable}, \texttt{stackable},
\texttt{reversible} --- are derived from those checks and cannot be declared.

Two limits belong with the definition. Patch identifiers use their own prefix
and stay inside the payload, so the record-level hash walk, which resolves
\texttt{sha256:} strings, does not see them; P-02 and P-03 check them instead,
and the record-level rule was left narrow rather than silently widened. And
\texttt{num\_changed} can only be recomputed when the base's shapes are known
and no operation names the same row twice; where it cannot be recomputed, the
check reports that the count is declared rather than assuming either answer.

\section{Record Qualifiers}

A record can carry evidence that no check asks about: that a baseline run
exists, that every run measured one population, that another operator rebuilt
the claim without disclosing a cost, that the invariants were re-measured and
not only the target metric. The status vocabulary is deliberately narrower than
that evidence, for the reason section 6 of the paper gives for status itself: a
value a record may declare is a value a record may claim. So the four
measurements are named as \emph{qualifiers}, derived by the test from the record
and never read from it.

\begin{table}[h]
\centering
\begin{tabular}{@{}ll@{}}
\toprule
Qualifier & State it names \\
\midrule
\texttt{baseline-present} & at least one run applies no intervention \\
\texttt{single-population} & every run reports the same unit count and sample count \\
\texttt{unpriced-reproduction} & another operator's run is uncosted and reports the claim \\
\texttt{invariants-replicated} & a counted run reports every invariant metric \\
\bottomrule
\end{tabular}
\end{table}

Each one says less than its name suggests, and the companion document states
what it does not mean. The one most likely to be misread is
\texttt{unpriced-reproduction}: it does not say the claim was replicated. It
names the run that the cost rule excludes from counting, and it makes that run
visible instead of letting a reader infer replication from a status.

\section{The adoption surface}

A format is adopted when the tools that produce the artifacts emit it. An
adapter for one such tool now reads the run directory that tool publishes --- a
manifest, an edit as named tensors, and a short run description --- and writes a
record and a patch. It delegates every judgement to the test: it does not
restate a rule, it writes a record and the test says whether the record
conforms. It does refuse one thing. MIR-1-04 fails a run that states no cost, so
the adapter will not write a record whose runs state none unless it is told to,
and then it writes a record that fails that check. An adapter that quietly
reproduced the one failure of the worked example would remove the reason the
example exists.

The adapter has been exercised against a fixture cut in the shape the upstream
tool publishes, not against the tool itself, and that is stated where the
adapter is documented. It is worth stating here too, because it is the same
point the paper makes about conformance: a record produced with the wrong
convention for an edit is still conforming, and the test cannot tell. The format
makes form checkable. It does not make the artifact true.

\section{Residual limits}

The limits of version 1.0 remain. Invariant bounds are declared, not evaluated:
the test checks that a bound is stated as a number and that the invariant set
reaches outside the claimed score, not that the bound was met. Operator identity
is taken from what the record declares. No third implementation exists, so
implementer independence is still unmet; the two implementations that do exist
agree byte for byte on every record, on every mutation in the repository, and on
a corpus of four hundred random documents, which is evidence that the text is
implementable and not evidence of independence.

Two limits are new. A patch is checked as a payload, so a patch that is
well-formed and applies to the wrong checkpoint is caught only if its declared
base disagrees with the record's base. And an addendum is still one account by
one collective author: it records what changed, and a reader who wants to know
whether the change was warranted should run the test rather than read this.

\section{Checking the claims in this addendum}

Every number above is reproducible from the release. The eleven verdicts for
any record come from the validator, which needs nothing installed:

\begin{verbatim}
cd MIR
python3 validate.py records/MIR-R1.json
python3 ci/check.py
\end{verbatim}

The first command ends in a failure of exactly one check, on cost disclosure,
which is the paper's subject. The second asserts every expected answer in the
release, including the ones that must fail, and runs both implementations
against each other. The patch examples are in the repository, each failing one
patch check in turn, so that the five verdicts are not merely reported but seen
to be reachable.

\end{document}
